\chapter{讨\hspace{6pt}论}

\section{与现有固井评价方法的关系}

井完整性研究强调套管和水泥屏障需要结合服役场景进行评价\cite{Kiran2017WellIntegrityReview,Bai2016WellIntegrityReview}。传统 CBL 和 VDL 通过套管波幅度、衰减和波形连续性评价胶结状态，其优势是技术成熟、解释路径清晰，但结果受套管、水泥、工具和地层条件共同影响\cite{Pardue1963CementBond,Wang2016AcousticCementMethods}。周向分区和方位阵列工具通过增加周向敏感单元改善局部缺陷识别\cite{Bigelow1990SegmentedBond,Lu2014AzimuthalAcoustic,Zuo2020AzimuthCement}。超声脉冲回波和周向成像则能提供更直观的套管后介质图像\cite{Froelich1982CementEvaluation,VanKuijk2005UltrasonicImager}。本文并不替代这些物理测量，而是研究在已有 XSI 和 CAST 资料之间建立可追溯的统计映射。

与直接方位成像方法不同，本文没有可靠的共同绝对方位零点。方位声学文献中，接收阵列的方向解释建立在工具几何和参考方向已知的条件上\cite{Che2014PhasedArcReception,Che2017AzimuthalReceiver}。本文主要实验井段整体接近直井，Inclination 中位数约为 $0.4705^\circ$，Relative Bearing 又存在频繁跳变；Inclination 仅描述井斜大小，不能替代绕工具轴的旋转角。本文因而不使用绝对方位位置作为监督，而以 FFT 幅值保留环向结构的频率组成，使同一结构在未知循环移位下具有相同标签。其价值是降低未知方位移位造成的监督不一致，而不是提高仪器方位分辨率或恢复窜槽的绝对方位。

\section{标签设计的价值与信息损失}

一维窜槽占比标签对方位零点完全不敏感，但把连续宽通道、多个离散异常和不同幅值分布压缩为相同占比。FFT 幅值标签保留周向总量和不同尺度的周期结构，因此比单一占比包含更多结构信息。离散傅里叶平移关系给出了其循环平移不变性的数学依据，傅里叶描述子研究也说明频域幅值可用于表达不依赖起点的形状特征\cite{Zahn1972FourierDescriptors,Persoon1977FourierDescriptors,Kuhl1982EllipticFourier}。

该设计同时存在不可逆损失。舍弃相位后，模型无法恢复窜槽位于哪个绝对方位；不同序列还可能具有相同幅值谱。只保留前 30 个系数又进一步削弱细尺度周向变化。因此，本文题目中的“结构特征反演”应理解为对旋转不变频谱特征的预测，不应解释为重构完整 CAST 图像或恢复绝对窜槽位置。

阈值 $2.5~\mathrm{MRayl}$ 是哈里伯顿公司在本项目 CAST 成像解释中采用的项目约定，不是本文提出或通过优化得到的普适行业标准。全井 $Z_c$ 中的 14 个负值缺乏合理物理意义，历史代码没有显式屏蔽；但位置复核表明它们全部位于本文目标井段之外，目标插值网格中没有负值，因而未进入一维占比、严重度或 FFT 标签以及模型训练、验证和测试。因此，这些负值不影响本文主要实验结果，而阈值选择、参考工具测量和频谱变换仍会共同影响派生标签效度。实验结果描述的是对该派生表征的学习程度，而不是对地下水泥环真实几何的直接测量精度。

\section{CWT 表征与网络骨干的作用}

CWT 将非平稳声波映射为局部时间--频率表示，较整体频谱保留到达时间。声波测井时频小波研究和水泥评价迁移学习应用均表明，这类表示适合处理波形中随时间变化的频率成分\cite{Saito1997WaveletAcoustic,Abdollahian2024TransferCement}。本文实验一也表明，在随机样本划分下，CWT 输入包含与 CAST 派生类别相关的可学习信息。

EfficientNetV2B0 在本文中只是特征提取骨干。其训练效率和结构缩放特点来自原始方法\cite{Tan2021EfficientNetV2}，不能作为本文网络创新。当前数据量、输出维度和早期过拟合现象表明，单纯增加网络容量未必改善深度区间迁移。与水泥评价分类、固井参数预测和声波日志重构研究相比\cite{Viggen2020AutomaticCement,Liu2025WideDeepCement,Wang2023ExplainableSonic}，本文的输入、标签和验证目标均不同，不能直接比较准确率或误差数值。

\section{单井深度区间隔离结果的解释}

实验三测试 $R^2=0.106634$，Pearson 为 0.351950，Spearman 为 0.480519。正相关和正 $R^2$ 表明模型在未见连续深度区间中保留了一部分与标签变化相关的信息，但解释方差有限。验证损失第 2 轮达到最佳，说明模型很快拟合主要模式，之后没有获得稳定的验证改善。

深度区间隔离减少了相邻样本直接混入不同集合的风险。结构化数据验证研究表明，分块验证通常比随机划分更能反映迁移到新空间或时间区间的误差\cite{Roberts2017CrossValidation,Ploton2020SpatialValidation}。但本文仍只有单井、单接收环数据。不同井之间可能存在套管规格、地层、井液、仪器增益、姿态和标签分布差异，单井深度隔离不能替代跨井验证，也不能据此判断工业部署条件下的稳定性。

\section{零值基线与稀疏标签}

实验三的零值基线 MAE 为 0.069751，低于模型的 0.079025；模型则在 RMSE 和 $R^2$ 上优于零值基线。该差异来自指标对误差分布的权重不同：MAE 对每个元素等权，大量零或近零标签使全零预测占优；RMSE 对少量大误差更敏感，模型对部分非零结构的预测可以降低平方误差。

因此，“模型是否优于基线”没有脱离指标的唯一答案。本文只能陈述：模型对非零结构显示了一定学习能力，但没有在平均绝对误差上击败最简单的稀疏标签基线。训练集中位数基线的 RMSE 与模型非常接近，也说明总体结果仍受标签中心分布强烈约束。后续可考虑非零区域指标、分严重度评价或与工程风险相关的代价函数，但在没有新增实验前不能把这些设想写成已验证改进。

\section{低频误差与高严重度低估}

低阶 FFT 系数承载周向总量和大尺度结构，其目标幅值较大，测试绝对误差也高。高阶系数 MAE 较小主要可能来自目标幅值较小，而不一定表示模型更容易学习细尺度结构。后续分析应同时报告归一化误差、目标能量和重构差异，避免只依据绝对 MAE 判断不同频率的难度。

测试样本中没有平均积分严重度超过 20 的样本，5--20 区间仅 75 个。较高严重度样本的预测峰值普遍不足，提示模型可能向常见低值收缩。样本不平衡、Huber 损失对大误差的线性增长、高维输出中的稀疏元素和区间分布变化都可能参与这一现象，但本研究没有设计消融实验区分这些因素。该问题直接限制了模型在高风险井段筛查中的用途。

\section{Grad-CAM 的解释边界}

Grad-CAM 依据指定目标输出对卷积特征图的梯度生成粗粒度响应区域\cite{Selvaraju2017GradCAM}。对于本文的多输出回归，不同深度位置、不同 FFT 系数或不同系数组合对应不同目标，热图必须随目标定义一起报告。热区只能说明该预测在局部线性梯度意义下对某些时频区域敏感，不能解释为网络参数，也不能证明真实物理因果。

现有候选热图只能部分追溯到早期 FFT 回归分析，无法同时确定每张图生成时的 checkpoint、训练运行、样本深度和数据划分，且其标量目标是全部回归输出之和。因此，本文没有把这些图用于实验三解释。后续若建立完整的可追溯分析，应先固定测试样本和目标系数，再比较不同误差样本、接收通道和声波到达时窗；即使完成这些步骤，结论仍应限定为模型响应与物理时窗之间的空间一致性，而非机理证明。

\section{研究有效性威胁}

本文的主要有效性威胁包括：CAST 的 $2.5~\mathrm{MRayl}$ 阈值仅适用于本项目数据约定；Relative Bearing 的正式参考方向和 CAST 零点不能形成与 XSI 唯一对应的共同方位基准；XSI 波形存在少量有符号 24 位端点削顶，历史预处理没有专门恢复；严格验证只覆盖单井第 3 接收环；较高严重度样本有限；输出幅值谱不能恢复绝对方位。XSI 端点数量极少，经项目数据核查，其对总体实验结果的影响可以忽略，但仍属于输入数据质量现象。CAST 全井 14 个负值因位于目标井段之外而没有进入本文标签，不构成主要实验指标的不确定来源。上述限制分别影响标签解释、输入质量、外部效度和结论范围。

上述问题并不否定当前实验的研究价值，但决定了结果应被表述为方位失配条件下的一次可追溯方法验证。只有在补充工具姿态元数据、跨井数据和更完整的标签质量控制后，才能进一步讨论绝对方位、跨井稳定性和工程应用。

\section{本章小结}

本章将本文方法置于传统声幅、方位阵列、超声成像和数据驱动固井评价的技术脉络中。FFT 幅值标签在不使用绝对方位位置的同时保留环向结构的频率组成，但以丢失相位、绝对方位和部分细尺度结构为代价；CWT 与 EfficientNetV2B0 提供了实现路径，而非结论本身。单井深度区间隔离结果显示有限可学习性，同时暴露了零值基线 MAE、低阶系数误差、高严重度低估和跨井证据缺失等限制。
